\section{\heiti 抽水蓄能财务模型特征分析}

本节对研究对象——"数字化系统财务模型边界\_抽水蓄能\_v12.xlsx"——进行系统性结构分析。该模型是某抽水蓄能电站项目的完整财务评价模型，涵盖从投资概算、资金筹措、折旧摊销、成本费用到利润、现金流、资产负债的全链条计算。分析揭示了模型的结构特征和关键业务复杂性，为后续方法设计提供依据。

\subsection{\heiti 模型整体结构}

该模型包含14个工作表，按计算逻辑的依赖关系可分为三个层次：输入层、计算层和输出层。输入层包含原始参数和时间序列数据，计算层执行中间变量推导，输出层生成最终财务指标。各层工作表如表~\ref{tab:sheets} 所示。

\clearpage

\begin{table*}[t]
	\centering
	\caption{模型工作表结构与公式分布}
	\label{tab:sheets}
	\vspace {-2.5mm}
	\begin{center}
		\begin{tabular}{lccccc}
			\toprule
			Sheet名称 & 层级 & 非空单元格 & 公式数 & 公式占比/\% & 跨Sheet引用数 \\
			\hline
			参数输入表 & 输入层 & 682 & 442 & 64.8 & 12 \\
			时间序列 & 输入层 & 3,218 & 1,843 & 57.3 & 428 \\
			投产\&达产比例 & 输入层 & 7,176 & 4,675 & 65.1 & 892 \\
			投资概算明细 & 输入层 & 562 & 306 & 54.4 & 45 \\
			表1-资金筹措及还本付息表 & 计算层 & 13,372 & 10,301 & 77.0 & 2,341 \\
			表2-折旧摊销表 & 计算层 & 7,559 & 6,765 & 89.5 & 1,876 \\
			表3-成本费用表 & 计算层 & 4,330 & 3,528 & 81.5 & 1,203 \\
			表4-收入税金表 & 计算层 & 2,788 & 2,253 & 80.8 & 892 \\
			表5-利润表-资本金 & 输出层 & 3,411 & 3,060 & 89.7 & 567 \\
			表6-现金流量表-资本金 & 输出层 & 2,055 & 1,802 & 87.7 & 389 \\
			表7-利润表-全投资 & 输出层 & 2,540 & 2,249 & 88.5 & 412 \\
			表8-现金流量表-全投资 & 输出层 & 1,610 & 1,424 & 88.4 & 298 \\
			表9-现金流量表-财务计划 & 输出层 & 2,117 & 1,829 & 86.4 & 345 \\
			表10-资产负债表 & 输出层 & 1,996 & 1,632 & 81.8 & 297 \\
			\hline
			合计 & & 57,577 & 42,109 & 73.1 & 9,199 \\
			\bottomrule
		\end{tabular}
	\end{center}
\end{table*}

模型整体公式占比73.1\%，说明该模型以自动化计算为主，而非静态数据汇总。其中计算层Sheet的公式占比最高（表2-折旧摊销表达89.5\%），反映了中间计算的密集程度。

\subsection{\heiti 关键结构特征发现}

通过对14个工作表的逐层分析，本文发现两个此前未被报道的业务复杂性特征。

\textbf{资本金/全投资双视角结构。} 该模型同时包含"资本金"和"全投资"两个视角的利润表和现金流量表（表5-表8），这在能源项目财务评价中具有特殊业务含义：资本金视角仅计算股东实际投入资金的回报率（资本金IRR），而全投资视角将全部建设资金（含借款）视为投资，计算项目整体盈利能力（全投资IRR）。两个视角在计算逻辑上共享部分输入（如收入、成本），但在资金筹措和利息处理上存在显著差异——资本金视角需扣除借款本金和利息支出，全投资视角则将利息内含。这种双视角结构导致模型中存在大量"同源异算"的计算路径：同一收入数据分别流入两套利润表和现金流表，增加了跨Sheet依赖的复杂度。现有电子表格分析方法尚未覆盖此类"一输入多输出"的业务结构。

\textbf{跨业务Sheet现象。} 进一步分析发现，部分Sheet的名称与其实际包含的计算逻辑存在语义冲突。典型例子是"时间序列"Sheet：从名称看应为纯时间维度数据，但实际内嵌了折旧摊销时间轴的计算逻辑，导致该Sheet的行标题同时包含"时间"关键词和"折旧"关键词。这一现象我们称之为"跨业务Sheet"——单一Sheet内承载了多个业务模块的计算逻辑，导致Sheet名称（上下文特征之一）与行标题（另一上下文特征）之间存在语义歧义。消融实验（第6.2节）表明，这是规则方法天然边界所在：28/35条上下文规则方法的错误集中于该Sheet。

\subsection{\heiti 公式类型分布}

对42,109个公式按语法结构进行分类统计，结果如表~\ref{tab:formula_dist} 所示。

\begin{table}[H]
	\centering
	\caption{公式类型分布}
	\label{tab:formula_dist}
	\vspace {-2.5mm}
	\begin{center}
		\begin{tabular}{lcc}
			\toprule
			公式类型 & 数量 & 占比/\% \\
			\hline
			算术运算 & 18,557 & 44.1 \\
			跨Sheet引用 & 11,668 & 27.7 \\
			聚合函数（SUM/AVERAGE等） & 5,546 & 13.2 \\
			条件判断（IF类） & 4,933 & 11.7 \\
			条件聚合（SUMIF/COUNTIF等） & 1,402 & 3.3 \\
			财务指标函数（IRR/NPV等） & 3 & 0.0 \\
			\bottomrule
		\end{tabular}
	\end{center}
\end{table}

算术运算和跨Sheet引用合计占比71.8\%，说明该模型以数值计算和跨表数据传递为主。聚合函数（13.2\%）和条件判断（11.7\%）的占比较低，但覆盖了融资、折旧、成本等关键业务逻辑。财务指标函数仅3个（IRR/NPV类），但作为模型的最终输出目标，其计算链覆盖全模型，是路径追溯的核心锚点。

\subsection{\heiti 依赖图（DAG）统计}

基于公式引用关系构建有向无环图（DAG），统计结果如表~\ref{tab:dag_stats} 所示。

\begin{table}[H]
	\centering
	\caption{依赖图（DAG）统计}
	\label{tab:dag_stats}
	\vspace {-2.5mm}
	\begin{center}
		\begin{tabular}{lc}
			\toprule
			统计项 & 数值 \\
			\hline
			总节点数 & 62,160 \\
			总边数 & 78,472 \\
			公式节点数 & 42,109 \\
			孤立节点数 & 11,004 \\
			平均入度 & 1.26 \\
			平均出度 & 1.26 \\
			最大入度 & 27 \\
			最大出度 & 1,021 \\
			跨Sheet边数 & 9,146 \\
			源节点数 & 25,629 \\
			汇节点数 & 24,746 \\
			最长路径长度 & 730 \\
			循环引用 & 无 \\
			\bottomrule
		\end{tabular}
	\end{center}
\end{table}

DAG验证结果为无环，说明该模型不存在循环引用。最大出度1,021的节点为某参数单元格，被1,021个公式共同引用，体现了参数敏感性的集中特征——修改该参数将影响1,021个下游计算结果。最长路径730步表明从某个输入参数到最终输出指标需经历730次计算传递，远超人工可追溯的范围。
